% ===== 5 实验结果与分析 =====
\section{实验结果与分析}
\label{sec:r-results-final}

\subsection{实验材料与评价范围}

本文实验材料由三张自然堆积真实图像及其端到端输出、竞赛层合成单元测试和原型运行阶段记录的算法时序组成。三张真实图像均缺少一一配对的机械筛分逐筛称量、人工实例 mask 和材料形貌标签，因此本章评价四类可由现有数据直接观察的结果：端到端处理是否完整运行，不同场景是否产生有区分度的级配描述，数量统计与质量代理统计如何改变分布，以及规则模块与报告实现是否保持一致。表\ref{tab:r-final-evidence}列出各类证据对应的结论范围。

\begin{table}[t]
\centering
\caption{本文使用的证据及其对应的结论范围。}
\label{tab:r-final-evidence}
\footnotesize
\renewcommand{\arraystretch}{1.12}
\begin{tabular}{>{\raggedright\arraybackslash}p{3.0cm}>{\raggedright\arraybackslash}p{4.4cm}>{\raggedright\arraybackslash}p{5.6cm}}
\toprule
证据 & 内容 & 可支持的结论 \\
\midrule
3 张自然堆积图像 & 原图、实例结果、逐颗粒几何、级配、形貌与异常输出 & 处理链可运行；不同场景形成不同的结构化描述 \\
数量/质量代理统计 & 同一实例集合分别按等权和 $d^3$ 加权 & 统计口径改变 $D_p$ 与级配解释 \\
27 项合成单元测试 & 筛分映射、加权分位数、形貌、异常、报告逻辑 & 数学定义与代码实现保持一致 \\
历史算法时序 & 分割、几何测量和筛分统计的阶段耗时 & 计算阶段可嵌入分钟级现场流程 \\
\bottomrule
\end{tabular}
\end{table}

\subsection{真实自然堆积图像的端到端输出}

三张样本分别为 \texttt{agg\_001}、\texttt{agg\_005} 和 \texttt{agg\_029}，图像尺寸均约为 3024$\times$4032，现有产物使用的尺度为 0.208~mm/px。三张图共检测到 3858 个实例，并生成逐颗粒几何表、筛分等效粒径、质量代理级配、投影形貌标签、异常候选和场景级可视化。表\ref{tab:r-final-demo-main}给出异常规则筛选后的正常实例结果，参数统一为 $\boldsymbol{\theta}=(1,0,0)$、$\gamma=3$。

\begin{table}[t]
\centering
\caption{三张自然堆积样本的视觉质量代理级配输出。}
\label{tab:r-final-demo-main}
\footnotesize
\renewcommand{\arraystretch}{1.12}
\begin{tabular}{lrrrrll}
\toprule
样本 & 实例数 & 正常实例 & $D_{10}$ & $D_{50}$ & $D_{90}$ & 主粒级 \\
 & & & \multicolumn{3}{c}{/ mm} & \\
\midrule
\texttt{agg\_001} & 1924 & 1193 & 11.1 & 23.6 & 30.6 & 25--31.5~mm（34.2\%） \\
\texttt{agg\_005} & 968 & 703 & 15.6 & 27.5 & 38.0 & 25--31.5~mm（29.7\%） \\
\texttt{agg\_029} & 966 & 520 & 6.7 & 14.6 & 30.6 & 5--10~mm（32.7\%） \\
\bottomrule
\end{tabular}
\end{table}

三个场景呈现出不同的粒径结构。\texttt{agg\_005} 的粗粒尾部最强，$D_{50}$ 和 $D_{90}$ 均为三组最高；\texttt{agg\_029} 的分布中心更靠细端，主粒级位于 5--10~mm。按当前质量代理统计，\texttt{agg\_001} 的 20--31.5~mm 合计约 55.3\%，\texttt{agg\_005} 的 25--40~mm 约 55.6\%，\texttt{agg\_029} 的 5--16~mm 约 51.9\%。统一参数下仍能保留这些场景差异，说明端到端链条没有把输入压缩成近似相同的输出分布。

\begin{figure}[t]
\centering
\begin{subfigure}{0.32\linewidth}
\centering
\includegraphics[width=\linewidth]{agg_001_IG2_full_set_cp_SAM_pred_grains_re_scaled_gsd_comparison.png}
\caption{\texttt{agg\_001}}
\end{subfigure}\hfill
\begin{subfigure}{0.32\linewidth}
\centering
\includegraphics[width=\linewidth]{agg_005_IG2_full_set_cp_SAM_pred_grains_re_scaled_gsd_comparison.png}
\caption{\texttt{agg\_005}}
\end{subfigure}\hfill
\begin{subfigure}{0.32\linewidth}
\centering
\includegraphics[width=\linewidth]{agg_029_IG2_full_set_cp_SAM_pred_grains_re_scaled_gsd_comparison.png}
\caption{\texttt{agg\_029}}
\end{subfigure}
\caption{三张样本的数量型与质量代理型粒径分布，用于比较场景结构和统计口径。}
\label{fig:r-final-three-gsd}
\end{figure}

\subsection{数量分布与质量代理分布}

表\ref{tab:r-final-number-mass}在完全相同的实例集合上比较等权统计和 $d^3$ 加权统计。三张图的数量型 $D_{50}$ 为 5.2--7.6~mm，质量代理 $D_{50}$ 则为 14.1--27.5~mm；三个场景的分布中心均向粗端移动。这一变化直接对应第\ref{sec:r-sieve-final}节提出的统计语义问题：实例数量多并不意味着对应粒级在质量上占比高。

\begin{table}[t]
\centering
\caption{数量型与质量代理型分位数比较。}
\label{tab:r-final-number-mass}
\footnotesize
\renewcommand{\arraystretch}{1.12}
\begin{tabular}{lrrrrrr}
\toprule
& \multicolumn{3}{c}{数量型 / mm} & \multicolumn{3}{c}{质量代理 / mm} \\
\cmidrule(lr){2-4}\cmidrule(lr){5-7}
样本 & $D_{10}$ & $D_{50}$ & $D_{90}$ & $D_{10}$ & $D_{50}$ & $D_{90}$ \\
\midrule
\texttt{agg\_001} & 3.1 & 6.0 & 17.3 & 10.6 & 23.5 & 30.6 \\
\texttt{agg\_005} & 3.3 & 7.6 & 24.5 & 15.4 & 27.5 & 38.0 \\
\texttt{agg\_029} & 3.0 & 5.2 & 9.2 & 5.9 & 14.1 & 30.6 \\
\bottomrule
\end{tabular}
\end{table}

自然堆积图像中，小颗粒和小目标在实例数量上占比较高，但单颗粒质量通常较低；少量粗颗粒的数量占比有限，却会在 $d^3$ 权重下获得更高贡献。\texttt{agg\_001} 提供了一个直接例子：$<5$~mm 实例约占全部检测数量的 38.0\%，将这些目标从正常口径中排除后，质量代理 $D_{50}$ 只由约 23.50~mm 变化到 23.56~mm。大量小实例因此会强烈影响数量分布，却未必同等影响质量代理分布。

同一机制也放大了粗尺度 merge 的风险。两个相邻颗粒若被合并为一个偏大实例，粒径误差会进一步通过 $d^3$ 权重进入累计分布，从而影响粗端级配和高分位数。实例分割在这里不仅决定颗粒数量，也决定后续质量代理中哪些对象会获得高权重。

这组对比支持的结论是数量分布与质量相关统计需要分开建模。$\gamma=3$ 在本文中仍是几何相似性先验，其作用由上述行为性对比展示，当前数据并未对质量指数本身进行机械称量校准。

\subsection{投影形貌与异常候选}

同一实例几何继续用于形貌和异常输出。表\ref{tab:r-final-shape}显示，三张图的普通类约占 55.0\%--64.1\%，针片状候选约占 14.0\%--20.9\%；\texttt{agg\_005} 的针片状候选比例相对较高，\texttt{agg\_029} 的圆形比例相对较高。这些类别由 AR、圆度和凸度阈值定义，用于描述当前视角下的二维轮廓组成。

\begin{table}[t]
\centering
\caption{三张样本的二维投影形貌规则输出（占全部检测实例比例，\%）。}
\label{tab:r-final-shape}
\footnotesize
\begin{tabular}{lrrrr}
\toprule
样本 & 针片状候选 & 圆形 & 棱角状 & 普通 \\
\midrule
\texttt{agg\_001} & 16.3 & 8.0 & 16.4 & 59.3 \\
\texttt{agg\_005} & 20.9 & 6.6 & 17.6 & 55.0 \\
\texttt{agg\_029} & 14.0 & 11.0 & 11.0 & 64.1 \\
\bottomrule
\end{tabular}
\end{table}

\begin{figure}[t]
\centering
\begin{subfigure}{0.24\linewidth}\centering
\includegraphics[width=\linewidth]{agg_001_IG2_full_set_cp_SAM_pred_grains_re_scaled_shape_needle_flaky.png}
\caption{针片状候选}
\end{subfigure}\hfill
\begin{subfigure}{0.24\linewidth}\centering
\includegraphics[width=\linewidth]{agg_001_IG2_full_set_cp_SAM_pred_grains_re_scaled_shape_round.png}
\caption{圆形}
\end{subfigure}\hfill
\begin{subfigure}{0.24\linewidth}\centering
\includegraphics[width=\linewidth]{agg_001_IG2_full_set_cp_SAM_pred_grains_re_scaled_shape_angular.png}
\caption{棱角状}
\end{subfigure}\hfill
\begin{subfigure}{0.24\linewidth}\centering
\includegraphics[width=\linewidth]{agg_001_IG2_full_set_cp_SAM_pred_grains_re_scaled_shape_regular.png}
\caption{普通}
\end{subfigure}
\caption{\texttt{agg\_001} 的二维投影形貌高亮，用于检查规则定位与几何定义。}
\label{fig:r-final-shape-montage}
\end{figure}

异常规则在当前三张样本中主要触发 $<5$~mm 小目标，$>50$~mm 和疑似泥团均为 0，如表\ref{tab:r-final-anomaly}所示。现有样本因此能够展示异常分支的规则执行和输出位置，但没有提供 oversized 或泥团阳性样本用于估计召回率。

\begin{table}[t]
\centering
\caption{三张样本的异常规则输出（实例数）。}
\label{tab:r-final-anomaly}
\footnotesize
\begin{tabular}{lrrrr}
\toprule
样本 & 正常 & $>50$~mm & $<5$~mm & 疑似泥团 \\
\midrule
\texttt{agg\_001} & 1193 & 0 & 731 & 0 \\
\texttt{agg\_005} & 703 & 0 & 265 & 0 \\
\texttt{agg\_029} & 520 & 0 & 446 & 0 \\
\bottomrule
\end{tabular}
\end{table}

\subsection{软件一致性与运行效率}

仓库中的竞赛层测试覆盖筛分等效粒径、质量权重、加权分位数、形貌规则、异常优先级和报告生成，共 27 项合成测试，当前记录全部通过。合成测试使用已知输入检查公式和规则对应的输出，因此评价的是数学定义与代码实现的一致性。现实骨料的误差还会受到二维投影、实例分割和质量代理假设影响，这部分不由单元测试给出。

原型历史运行记录约为实例分割 17~s、几何测量 9~s、筛分与报告统计小于 1~s，即算法阶段约 27~s/张。完整现场流程还包含尺度确认、图像质检、首次模型加载、结果复核以及必要的重拍或复跑。现有时序说明计算阶段本身只占 30~min 时间窗口的一小部分，第\ref{sec:r-deploy-final}节据此安排现场操作预算。

\subsection{结果小结}

现有证据支持三个主要观察。首先，三张真实自然堆积图像均完成了从实例分割、毫米几何到级配、投影形貌和异常候选的端到端处理。其次，统一参数下三个场景保持不同的 $D_p$、主粒级和分布形状，其中数量型与质量代理型 $D_{50}$ 的系统性差异直接显示了统计口径对级配解释的影响。最后，27 项合成测试和现有时序记录分别给出了软件一致性与计算成本方面的证据。

本文当前评价范围不包含标准机械筛分绝对误差、实例分割外部 P/R/F1，以及三维针片状或泥团材料识别准确率。全文对这些输出分别使用“视觉质量代理级配”“投影形貌”和“异常候选”等术语，以保持结论与现有证据一致。